\documentclass[11pt]{article}
\usepackage{tcee}
\addbibresource{3B40.bib}
\tcetema[Cohesión, política social y convergencia en la UE]{3.B.40}{La Cohesión Económica y Social en la Unión Europea: política regional e instrumentos presupuestarios y financieros. Política social y de empleo. El proceso de convergencia real en la Unión Europea.}
\tcefecha{14/11/2026}
\begin{document}
\tceetitulo


\section*{Introducción}
\minutos{2,5}

\begin{itemize}
\item \textbf{\emph{Enganche}}:
  \begin{itemize}
  \item «Dios da pan a quien no tiene dientes», dice el refrán: los recursos llegan a menudo a quien menos partido puede sacarles. Es una crítica que se oye a veces a la política de cohesión, que concentra sus fondos en las regiones menos desarrolladas, las que, en principio, menos capacidad tienen para aprovecharlos.
  \item Los datos la apoyan solo a medias. Desde 2000, Polonia ha pasado del 48~\% al 81~\% de la renta media de la Unión por habitante y Rumanía, del 26~\% al 78~\%; en el mismo tiempo, España ha bajado del 97~\% al 92~\%, Italia del 122~\% al 96~\% y Grecia del 85~\% al 68~\%, después de recibir durante décadas los mayores fondos de cohesión\nocite{Eurostat2026PPA}.
  \item La pregunta que recorre la exposición es: \emph{¿sirve la política de cohesión para que las regiones pobres de la Unión alcancen a las ricas, o solo compensa a quien pierde con la integración?}
  \end{itemize}

\item \textbf{\emph{Relevancia}}:
  \begin{itemize}
  \item \emph{Es la mayor política de inversión de la Unión}: unos 392.000 millones de euros en 2021-2027, casi 1/3 del marco financiero plurianual\nocite{CE2026Cohesiondata} \vertema{3.B.37}. En varios Estados del Este financia más de la mitad de la inversión pública.
  \item \emph{Está en plena discusión}: la Comisión propone que, desde 2028, la cohesión y la PAC se repartan en un solo plan nacional y regional por Estado, lo que las regiones y el Parlamento ven como una centralización\nocite{CE2025PropuestaFondo}.
  \item \emph{Para España}: ha sido el mayor beneficiario de los fondos estructurales hasta 2006 y recibe unos 36.700 millones en 2021-2027; aun así, su renta por habitante no ha convergido con la media desde 2008\nocite{CE2018Panorama64,CE2026Cohesiondata,Eurostat2026PPA} \vertema{4.A.17} \vertema{4.A.30}.
  \end{itemize}

\item \textbf{\emph{Contextualización}}:
  \begin{itemize}
  \item \emph{Qué es}: la \concepto{cohesión económica, social y territorial} es el objetivo de reducir las diferencias de desarrollo entre las regiones y el retraso de las menos favorecidas (art. 174 TFUE). Es una \emph{competencia compartida} (art. 4 TFUE) y se persigue sobre todo con fondos.
  \item \emph{Tiene 2 brazos}:
    \begin{itemize}
    \item la \concepto{política regional}, con el Fondo Europeo de Desarrollo Regional (FEDER), el Fondo de Cohesión y el Fondo de Transición Justa, que invierte en las regiones;
    \item la \concepto{política social y de empleo}, con el Fondo Social Europeo Plus (FSE+), normas mínimas y coordinación de las políticas nacionales, que actúa sobre las personas.
    \end{itemize}
  \item \emph{Los 2 se distinguen por su base jurídica}: la cohesión es competencia compartida y se hace con dinero; en lo social, la Unión solo fija mínimos en algunos ámbitos y en el empleo solo coordina (art. 5 TFUE). Por eso se estudian por separado.
  \item \emph{La convergencia real} es el acercamiento de las rentas por habitante (o de la productividad y el empleo) entre países y regiones; la \emph{nominal} es la de precios, tipos y finanzas públicas que pide Maastricht \vertema{3.B.43}.
  \end{itemize}

\item \textbf{\emph{Problemática}}:
  \begin{itemize}
  \item ¿Por qué necesita una política regional un mercado único que, según la teoría neoclásica, ya debería acercar las rentas?
  \item ¿Con qué instrumentos actúa la Unión, y cómo se gestionan?
  \item ¿Qué puede hacer la Unión en lo social y en el empleo, si los Estados se reservan los salarios y la seguridad social?
  \item ¿Ha habido convergencia real en la Unión, y qué parte se debe a los fondos?
  \end{itemize}

\item \textbf{\emph{Estructura}}:
  \begin{itemize}
  \item La exposición va de la política que más dinero mueve a la que más normas tiene, y termina con la prueba de los resultados:
    \begin{itemize}
    \item \concepto{I. Política regional} (\ref{sec:regional}): su relevancia, por qué existe (la integración tiende a concentrar la actividad), sus competencias, su marco jurídico y cómo se decide, su historia, sus instrumentos actuales, cómo se gestionan y sus retos.
    \item \concepto{II. Políticas social y de empleo} (\ref{sec:social}): sus bases, el Pilar Europeo de Derechos Sociales y sus instrumentos.
    \item \concepto{III. El proceso de convergencia real} (\ref{sec:convergencia}): el modelo de convergencia $\beta$ y $\sigma$ aplicado a las regiones de la Unión, la evidencia entre países y regiones, y el efecto de los fondos, con el modelo de regresión en discontinuidad.
    \end{itemize}
  \item Las 2 políticas se cuentan con el mismo orden: relevancia, justificación (teoría, competencias, marco jurídico y toma de decisiones), evolución, articulación y retos. La figura~\ref{figu:mapa-tema} resume el hilo.
  \end{itemize}
\end{itemize}


\begin{esquema}[unbreakable]
\textbf{3.B.40. La cohesión económica y social en la UE. Política social y de empleo. La convergencia real}\\[2pt]
Introducción\esqtiempo{2,5}\\
I. Política regional\esqtiempo{11}\\
\hspace*{1.5em}1. Importancia y justificación\\
\hspace*{1.5em}2. Marco jurídico\\
\hspace*{1.5em}3. Antecedentes: evolución histórica\\
\hspace*{1.5em}4. Instrumentos presupuestarios y financieros\\
\hspace*{1.5em}5. Los fondos regionales: cómo se gestionan\\
II. Políticas social y de empleo\esqtiempo{7}\\
\hspace*{1.5em}1. Política social\\
\hspace*{1.5em}2. Política de empleo\\
III. El proceso de convergencia real en la Unión Europea\esqtiempo{8,5}\\
\hspace*{1.5em}1. Caracterización teórica de la convergencia\\
\hspace*{1.5em}2. Evidencia sobre la convergencia en la Unión Europea\\
\hspace*{1.5em}3. Evidencia sobre el papel de la política regional\\
\hspace*{1.5em}4. Dificultades de medición\\
Conclusión\esqtiempo{1}
\end{esquema}


\begin{anotaciones}
\textbf{Mapa de la exposición}: cada bloque responde a una pregunta y deja abierta la del siguiente; a la derecha, lo que se cuenta en cada subbloque.
\begin{infografiacaja}
\caption{Mapa del tema: la pregunta de cada bloque, sus subbloques y el hilo que los une}\label{figu:mapa-tema}
\grafico[\linewidth]{mapa-tema}
\fuente{elaboración propia.}
\end{infografiacaja}
\end{anotaciones}

\begin{anotaciones}
\textbf{Notación y siglas.} El modelo de la integración y la localización tiene su notación en el anexo~\ref{anexo:capital-movil}. En el de convergencia, $y_t$ es la renta por habitante, $\gamma^{y}_{t,t-1}$ su crecimiento, $\beta_0$ y $\beta_1$ coeficientes, $\vect{z}_{t-1}$ los determinantes del estado estacionario, $\lambda$ la velocidad de convergencia y $\sigma_{y,t}^2$ la varianza del logaritmo de la renta. En el de la regresión en discontinuidad, $x_i$ es la renta por habitante de la región $i$ en \% de la media, $D_i$ el tratamiento (ser región del antiguo objetivo 1), $g_i$ su crecimiento y $\tau$ el efecto. FEDER: Fondo Europeo de Desarrollo Regional; FSE+: Fondo Social Europeo Plus; FTJ: Fondo de Transición Justa; RDC: Reglamento de disposiciones comunes; NUTS: nomenclatura de unidades territoriales estadísticas; PPA: paridad de poder adquisitivo; AROPE: población en riesgo de pobreza o exclusión social (del inglés \emph{at risk of poverty or social exclusion}).
\end{anotaciones}

% ==========================================================================
\section{Política regional}\label{sec:regional}
\minutos{11}

\begin{itemize}
\item Empezamos por la política regional, que es la que mueve casi todo el dinero de la cohesión.
\item \emph{Conclusión anticipada}: la integración no reparte por igual sus ganancias. Al abaratar el comercio, puede concentrar la actividad en el centro, y la política regional nació como compensación política y como corrección de esa tendencia. Hoy es una política de inversión de 7 años, con 5 fondos y un reglamento común, gestionada con los Estados y las regiones, cada vez más ligada a otras prioridades (la competitividad, la defensa, la vivienda) y con un futuro en discusión.
\end{itemize}

% --------------------------------------------------------------------------
\subsection{Importancia y justificación}\label{sec:justificacion}

\subsubsection{Importancia}

\begin{itemize}
\item \emph{En el presupuesto}: la cohesión es, con la PAC, una de las 2 grandes políticas de gasto. Desde 1994 se lleva en torno a 1/3 del presupuesto de la Unión, unos 0,3-0,4 puntos del PIB de la Unión al año \vertema{3.B.37}.
\item \emph{En los receptores}: para los Estados del Este, Portugal y Grecia, los fondos llegan a más del 2~\% de su PIB al año en los años fuertes de ejecución y a buena parte de su inversión pública\nocite{CE2024Informe9}.
\item \emph{En el mapa}: en 2024, 83 de las 244 regiones NUTS 2 de la Unión estaban por debajo del 75~\% de la renta media por habitante, el umbral de las regiones menos desarrolladas (gráfico~\ref{fig:mapa-pib})\nocite{Eurostat2026Regional}.
\end{itemize}

\subsubsection{Justificación política}

\begin{itemize}
\item \emph{Compensar a quien teme perder con la integración}: cada gran paso de la integración ha venido con más cohesión, como pago lateral para lograr la unanimidad:
  \begin{itemize}
  \item el FEDER (1975), tras la entrada del Reino Unido, Irlanda y Dinamarca, en parte para compensar al Reino Unido, que apenas recibía de la PAC;
  \item la duplicación de los fondos del Paquete Delors I (1988), con el mercado único; el \concepto{Informe Padoa-Schioppa} (1987) advirtió de que el mercado único podía agrandar las diferencias regionales si no se acompañaba de medidas de cohesión\nocite{PadoaSchioppa1987};
  \item el Fondo de Cohesión (1993), con el camino al euro, que España negoció en Maastricht y en el Consejo Europeo de Edimburgo (1992).
  \end{itemize}
\item \emph{La legitimidad}: los fondos hacen visible la Unión en cada región («cofinanciado por la Unión Europea») y ayudan a que los ciudadanos acepten la integración. Las regiones que se sienten olvidadas votan más a partidos contrarios a la Unión: es la «venganza de los lugares que no importan»\nocite{RodriguezPose2018}.
\item \emph{La solidaridad}: la cohesión es uno de los objetivos de la Unión (art. 3 TUE), que se piensa como una comunidad y no solo como un mercado.
\end{itemize}

\subsubsection{Justificación económica}

\begin{itemize}
\item \emph{La teoría neoclásica} promete convergencia sin política: con rendimientos decrecientes, el capital va a donde es escaso y rinde más, y el comercio iguala los precios de los factores. Las regiones pobres crecerían más deprisa \vertema{3.A.43}.
\item \emph{Pero hay varias razones para que no ocurra}, cada una con su remedio:
  \begin{itemize}
  \item \emph{fallos de mercado}: los mercados de capital imperfectos racionan el crédito a las regiones pobres; las infraestructuras y la formación son bienes públicos o tienen externalidades; puede haber fallos de coordinación, en los que ninguna empresa invierte si no invierten las demás (el gran impulso, \emph{big push}) \vertema{3.B.28};
  \item \emph{rendimientos crecientes y crecimiento endógeno}: si la tecnología y el capital humano se refuerzan a sí mismos y sus efectos desbordamiento son locales, las regiones ricas pueden crecer más \vertema{3.A.44};
  \item \emph{la causación circular acumulativa} de \autor{Gunnar Myrdal} (1957): el mercado tiende a concentrar la inversión donde ya hay actividad, y los «efectos de expulsión» sobre las regiones atrasadas pesan más que los de «difusión»\footnote{\autor{Myrdal} recibió el Premio Nobel de Economía en 1974, con \autor{Friedrich Hayek}, «por sus trabajos pioneros en la teoría del dinero y de las fluctuaciones económicas y por su penetrante análisis de la interdependencia de los fenómenos económicos, sociales e institucionales»\nocite{Nobel1974,Myrdal1957}.};
  \item \emph{la unión monetaria}: si las estructuras económicas son muy distintas, las perturbaciones son asimétricas y la política monetaria común no sirve a todos; acercar las estructuras y tener transferencias entre regiones ayuda a que funcione la zona del euro \vertema{3.B.16}.
  \end{itemize}
\item La razón más propia de la Unión es la que da la nueva geografía económica: el mercado único abarata el comercio, y abaratar el comercio cambia \emph{dónde} se produce.
\end{itemize}

\subsubsection{Integración, localización y el dilema de la cohesión}\label{sec:modelo-localizacion}

\begin{itemize}
\item \emph{La idea}: la \concepto{nueva geografía económica} de \textcite{Krugman1991Increasing} explica cómo una pequeña asimetría puede volverse una gran diferencia entre regiones \vertema{3.B.6} \vertema{3.B.10}. En su versión de \emph{capital móvil} (\emph{footloose capital}) de \textcite{MartinRogers1995}, el mercado grande atrae más que proporcionalmente a la industria, y tanto más cuanto más barato es comerciar\nocite{BaldwinEtAl2003Economic}.
\item \emph{Lo que enseña a la cohesión}:
  \begin{itemize}
  \item la integración no es neutral en el espacio: al abaratar el comercio, el mercado único puede llevar industria de la periferia al centro;
  \item las transferencias de renta a la periferia frenan esa concentración;
  \item no toda infraestructura ayuda a la periferia: la que la une con el centro puede vaciarla (la «carretera de doble sentido»), y la que abarata el comercio dentro de ella la refuerza;
  \item hay un dilema entre crecimiento y equidad espacial: concentrar la innovación en el centro hace crecer más a la Unión, pero agranda las diferencias \parencite{Martin1999}.
  \end{itemize}
\item El modelo completo, con sus supuestos, su equilibrio, el gráfico del mercado doméstico amplificado y su evidencia, está en el anexo~\ref{anexo:capital-movil}. Veamos ahora qué base jurídica tiene la Unión para actuar.
\end{itemize}

% --------------------------------------------------------------------------
\subsection{Marco jurídico}\label{sec:marco-regional}

\begin{itemize}
\item \emph{Delimitación de competencias}: la cohesión económica, social y territorial es una \concepto{competencia compartida} (art. 4.2.c TFUE): la Unión y los Estados pueden legislar, y los Estados lo hacen en la medida en que la Unión no lo haya hecho. Es distinta de las competencias exclusivas (art. 3: unión aduanera, competencia, política monetaria del euro, pesca y política comercial), de la coordinación de las políticas económicas y de empleo (art. 5) y de las de apoyo (art. 6: industria, educación, turismo, salud)\nocite{UE2016TFUE} (anexo~\ref{anexo:competencias}) \vertema{3.B.36}.
\item \emph{Marco jurídico} (título XVIII del TFUE, arts. 174-178):
  \begin{itemize}
  \item \emph{art. 174}: la Unión promueve un desarrollo armonioso y reduce las diferencias entre regiones y el retraso de las menos favorecidas, con especial atención a las zonas rurales, a las afectadas por una transición industrial y a las que tienen desventajas naturales o demográficas graves y permanentes (las islas, las de montaña, las muy poco pobladas y las transfronterizas). La dimensión \emph{territorial} se añadió con Lisboa (2009);
  \item \emph{art. 175}: los Estados coordinan sus políticas económicas para alcanzar estos objetivos, y la Unión los apoya con los fondos estructurales, el BEI y los demás instrumentos; la Comisión informa cada 3 años (el \emph{informe de cohesión});
  \item \emph{art. 176}: el FEDER corrige los principales desequilibrios regionales, con el desarrollo y el ajuste de las regiones atrasadas y la reconversión de las industriales en declive;
  \item \emph{art. 177}: los reglamentos de los fondos fijan sus objetivos, su organización y las normas comunes; el Fondo de Cohesión financia medio ambiente y redes transeuropeas de transporte;
  \item \emph{art. 178}: los reglamentos de aplicación del FEDER.
  \end{itemize}
\item \emph{Toma de decisiones}: los reglamentos de los fondos se aprueban por el \emph{procedimiento legislativo ordinario}, previa consulta al \concepto{Comité Económico y Social Europeo} (empresarios, sindicatos y otros grupos) y al \concepto{Comité de las Regiones} (cargos electos locales y regionales, creado en Maastricht y en funcionamiento desde 1994). El dinero de cada periodo, en cambio, lo fija el marco financiero plurianual, por unanimidad.
\end{itemize}

% --------------------------------------------------------------------------
\subsection{Antecedentes: evolución histórica}\label{sec:historia}

\begin{itemize}
\item La política regional ha crecido a saltos, con cada ampliación y cada paso de la integración (figura~\ref{figu:historia}):
  \begin{itemize}
  \item \emph{1957-1974: sin política regional}. El Tratado de Roma habla de un desarrollo armonioso y crea el Fondo Social Europeo, pero 6 Estados sin grandes diferencias (salvo el sur de Italia) no la necesitaban; el BEI financiaba proyectos en las regiones atrasadas.
  \item \emph{1975-1988: el FEDER por cuotas}. Tras la ampliación de 1973, el FEDER (1975) reparte dinero por cuotas nacionales para proyectos que eligen los Estados: es más una transferencia entre Estados que una política europea.
  \item \emph{1988-1999: la política de cohesión de verdad}. Con las entradas de Grecia (1981), España y Portugal (1986), cuya renta por habitante estaba entre el 50~\% y el 75~\% de la media, el Acta Única (1986) hace de la cohesión un objetivo del Tratado. La reforma de 1988 (Paquete Delors I) duplica los fondos y fija sus 4 principios: \emph{concentración} en las regiones más pobres (el objetivo 1, por debajo del 75~\%), \emph{programación} plurianual, \emph{asociación} con las regiones y los agentes sociales y \emph{adicionalidad} (los fondos se suman a la inversión nacional, no la sustituyen). En 1993 nace el Fondo de Cohesión, para los Estados que preparan el euro con menos renta.
  \item \emph{2000-2013: la ampliación al Este}. La Agenda 2000 prepara la entrada de 10 Estados en 2004 (y 2 en 2007), con un 20~\% más de población y solo un 5~\% más de PIB. Los fondos se orientan a la Estrategia de Lisboa (innovación) y la cohesión llega a superar a la PAC en el presupuesto\nocite{CE2018Panorama64}.
  \item \emph{2014-2020: resultados y condiciones}. Concentración temática en 11 objetivos, condiciones \emph{ex ante} y macroeconómicas, ligadas al Semestre Europeo\nocite{PE2026FichaCohesion} \vertema{3.B.45}.
  \item \emph{Desde 2020: más funciones con menos peso relativo}. La cohesión responde a la pandemia (REACT-EU), a la transición climática (Fondo de Transición Justa) y, en la revisión intermedia de 2025, a la defensa y la vivienda\nocite{UE2025RevisionIntermedia}.
  \end{itemize}
\end{itemize}

\begin{infografia}[H]
\centering
\caption{La política de cohesión en el tiempo: hitos y dotación de cada periodo}\label{figu:historia}
\grafico[0.97\textwidth]{historia-cohesion}
\fuente{elaboración propia a partir de la Comisión Europea\nocite{CE2018Panorama64,CE2026Cohesiondata} y de la propuesta de marco financiero 2028-2034. Dotación en miles de millones de euros corrientes de los fondos estructurales y de cohesión; 2000-2006, con los 21.700 millones para los nuevos Estados.}
\end{infografia}

\begin{itemize}
\item \emph{España en esta historia}: fue el mayor beneficiario de los fondos de 1989 a 2006 (14.200, 42.400 y 56.300 millones en los 3 periodos), con autovías, regadíos, depuradoras y formación; desde 2007, Polonia ocupa ese lugar\nocite{CE2018Panorama64}.
\end{itemize}

% --------------------------------------------------------------------------
\subsection{Instrumentos presupuestarios y financieros}\label{sec:instrumentos}

\subsubsection{El marco del periodo 2021-2027}\label{sec:marco-2021}

\begin{itemize}
\item \emph{Un reglamento para todos}: el \concepto{Reglamento de disposiciones comunes} (RDC, Reglamento 2021/1060) fija las reglas de 8 fondos de gestión compartida (los 4 de la cohesión, más el de pesca y 3 de interior), y cada fondo tiene además el suyo (FEDER y Fondo de Cohesión, 2021/1058; FSE+, 2021/1057; FTJ, 2021/1056; Interreg, 2021/1059)\nocite{UE2021RDC} (figura~\ref{figu:marco}).
\item \emph{5 objetivos políticos} (art. 5 RDC), que sustituyen a los 11 objetivos temáticos de 2014-2020: una Europa más competitiva e inteligente, más verde y baja en carbono, más conectada, más social e inclusiva y más próxima a los ciudadanos; el FTJ tiene un objetivo propio, ayudar a las regiones y personas más afectadas por la transición hacia la neutralidad climática\nocite{UE2021RDC}.
\item \emph{Los fondos} (importes de la Unión en euros corrientes)\nocite{CE2026Cohesiondata,PE2026FichaCohesion}:
  \begin{itemize}
  \item el \concepto{FEDER} (unos 214.000 millones, más 9.400 de Interreg): innovación, pymes, digitalización, energía, medio ambiente, transporte y desarrollo urbano; al menos el 8~\% para el desarrollo urbano sostenible\nocite{UE2021RDC};
  \item el \concepto{Fondo de Cohesión} (unos 39.000 millones, más 11.300 transferidos al Mecanismo «Conectar Europa»): medio ambiente y redes transeuropeas de transporte en los 15 Estados con renta nacional bruta por habitante inferior al 90~\% de la media (España no lo recibe desde 2014)\nocite{UE2021Categorias};
  \item el \concepto{FSE+} (unos 95.000 millones): empleo, educación e inclusión\nocite{UE2021FSE} (\ref{sec:politica-social});
  \item el \concepto{Fondo de Transición Justa} (FTJ; unos 22.000 millones iniciales, 19.700 hoy), nacido con el Pacto Verde para las regiones del carbón y la industria intensiva en emisiones, con dinero del marco y del NextGenerationEU\nocite{CE2026Cohesiondata};
  \item la \concepto{cooperación territorial europea} (Interreg): programas transfronterizos, transnacionales e interregionales.
  \end{itemize}
\item \emph{Las 3 categorías de regiones} (art. 108 RDC), según el PIB por habitante en PPA de 2015-2017 frente a la media de la UE-27, con su cofinanciación máxima de la Unión (art. 112; mapa en el gráfico~\ref{fig:mapa-categorias})\nocite{UE2021RDC}:
  \begin{itemize}
  \item \emph{menos desarrolladas} (menos del 75~\%): hasta el 85~\% del gasto. En España, Andalucía, Castilla-La Mancha, Extremadura, Ceuta y Melilla;
  \item \emph{en transición} (75-100~\%; la categoría pasó del 90~\% al 100~\% en 2021): hasta el 60~\% (70~\% si antes eran menos desarrolladas)\nocite{UE2021Categorias}. En España, 9: Galicia, Asturias, Cantabria, La Rioja, Castilla y León, la Comunidad Valenciana, Baleares, Murcia y Canarias;
  \item \emph{más desarrolladas} (más del 100~\%): hasta el 40~\% (50~\% si antes estaban por debajo). En España, el País Vasco, Navarra, Aragón, Madrid y Cataluña\nocite{UE2021Categorias}.
  \end{itemize}
\item \emph{El reparto entre Estados} sigue el \emph{método de Berlín}: la dotación de cada región depende de su población y de su distancia a la renta media, con primas por paro, paro juvenil, bajo nivel educativo, emisiones y migración. El 61~\% del dinero va a las regiones menos desarrolladas\nocite{UE2021RDC}.
\item \emph{Las reglas de calidad del gasto}:
  \begin{itemize}
  \item \emph{concentración temática}: los Estados más ricos deben dedicar al menos el 85~\% de su FEDER a los objetivos 1 y 2; los intermedios, el 40~\% al 1; los más pobres, el 25~\% al 1; todos, al menos el 30~\% al 2 (Reglamento 2021/1058, art. 4)\nocite{UE2021RDC}; en el FSE+, las reservas de \ref{sec:politica-social};
  \item \emph{el clima}: el 30~\% del FEDER y el 37~\% del Fondo de Cohesión;
  \item \emph{condiciones favorecedoras} (art. 15 RDC): 4 horizontales (contratación pública, ayudas de Estado, Carta de los Derechos Fundamentales, discapacidad) y 16 temáticas (una estrategia de especialización inteligente, planes de residuos, de transporte\ldots); si no se cumplen, ese gasto no se reembolsa;
  \item \emph{liberación automática de compromisos}: lo que no se gasta en 3 años ($n+3$) se pierde.
  \end{itemize}
\end{itemize}

\begin{infografia}[H]
\centering
\caption{La arquitectura de la política de cohesión 2021-2027}\label{figu:marco}
\grafico[0.97\textwidth]{marco-cohesion}
\fuente{elaboración propia a partir del RDC\nocite{UE2021RDC} y de la Comisión Europea\nocite{CE2026Cohesiondata} (importes de la Unión en la primera adopción de los programas).}
\end{infografia}

\begin{anotaciones}
\textbf{El mapa de las categorías}: las regiones menos desarrolladas forman un arco que va del Este (casi todo, salvo las capitales y algunas regiones más ricas) al sur de Italia, Grecia, el sur de España y Portugal y las regiones ultraperiféricas; las más desarrolladas, un eje que baja de Escandinavia, Dinamarca y los Países Bajos por el oeste y el sur de Alemania y Austria hasta el norte de Italia; en España, Madrid, Cataluña, Aragón, Navarra y el País Vasco.
\begin{figuracaja}
\caption{Las categorías de regiones de la política de cohesión 2021-2027}\label{fig:mapa-categorias}
\grafico[0.8\linewidth]{mapa-categorias}
\fuente{elaboración propia a partir de la Decisión de Ejecución (UE) 2021/1130\nocite{UE2021Categorias} (regiones NUTS 2 de 2021). Contornos: GISCO, Eurostat (© EuroGeographics para los límites administrativos)\nocite{GISCO2026NUTS}.}
\end{figuracaja}
\end{anotaciones}

\begin{itemize}
\item \emph{Quién recibe}: Polonia es, de lejos, el primer receptor en importe (gráfico~\ref{fig:cohesion-estados}, a la izquierda); por habitante, los bálticos, Eslovaquia, Croacia y Hungría (a la derecha). España es el tercero en importe, con 35.600 millones en los programas vigentes, pero recibe unos 724 euros por habitante, menos que la media de la Unión (815)\nocite{CE2026Cohesiondata}.
\item \emph{El Acuerdo de Asociación de España} (aprobado en noviembre de 2022): unos 36.700 millones de euros de la Unión, con 23.400 del FEDER, 11.300 del FSE+, 870 del FTJ y 1.120 del fondo de pesca; casi la mitad del FEDER va a las regiones menos desarrolladas\nocite{CE2026Cohesiondata,Hacienda2026FondosEuropeos} \vertema{4.A.17} \vertema{4.A.30}.
\end{itemize}

\begin{figure}[H]
\centering
\caption{La política de cohesión 2021-2027 por Estado, en miles de millones de euros y en euros por habitante}\label{fig:cohesion-estados}
\grafico[0.92\textwidth]{cohesion-estados}
\fuente{elaboración propia a partir de la Comisión Europea\nocite{CE2026Cohesiondata} (FEDER, Fondo de Cohesión, FSE+ y FTJ, sin Interreg; versión vigente de los programas, noviembre de 2026) y de Eurostat (población a 1 de enero de 2025). Media de la UE-27: el importe total entre la población.}
\end{figure}

\subsubsection{La revisión intermedia de 2025 y la STEP}\label{sec:revision-intermedia}

\begin{itemize}
\item \emph{La Plataforma de Tecnologías Estratégicas para Europa} (STEP, Reglamento 2024/795) permitió en 2024 reorientar fondos de cohesión hacia tecnologías digitales, limpias y biotecnológicas, con cofinanciación de hasta el 100~\% \vertema{3.B.37}. Hay unos 15.200 millones de la cohesión en prioridades STEP\nocite{CE2026Cohesiondata}.
\item \emph{La revisión intermedia} (propuesta de abril de 2025; Reglamentos 2025/1913 y 2025/1914, de septiembre de 2025) abrió los programas a nuevas prioridades\nocite{UE2025RevisionIntermedia}:
  \begin{itemize}
  \item la \emph{competitividad} y la descarbonización, la \emph{defensa} y la industria de doble uso (incluida la movilidad militar), la \emph{vivienda asequible}, la \emph{resiliencia hídrica}, la transición energética y las regiones de la frontera oriental;
  \item con incentivos: una prefinanciación extra del 4,5~\% (9,5~\% en la frontera con Rusia, Bielorrusia y Ucrania) si se reasigna al menos el 10~\% del programa, 10 puntos más de cofinanciación y un año más para gastar, hasta 2030.
  \end{itemize}
\item \emph{El resultado}: 25 Estados reasignaron 34.600 millones (casi el 10~\% de la cohesión) en 186 programas: 15.200 a la competitividad, 11.900 a la defensa, 3.300 a la vivienda, 3.100 al agua y 1.200 a la energía. España reasignó 3.200 millones, 2/3 a la competitividad\nocite{CE2026RevisionIntermedia}.
\item \emph{El debate}: la cohesión se ha vuelto la «caja» de la que salen las nuevas prioridades de la Unión. Para unos es flexibilidad; para otros, alejarla de su objetivo, la convergencia de las regiones atrasadas.
\end{itemize}

\subsubsection{Otros instrumentos}\label{sec:otros-instrumentos}

\begin{itemize}
\item \emph{Fondos de otras políticas con fines parecidos}: el FEADER, segundo pilar de la PAC, financia el desarrollo rural \vertema{3.B.38}; el fondo de pesca (FEMPA), las zonas costeras; el NextGenerationEU y su Mecanismo de Recuperación y Resiliencia, desde 2021, han sido en varios Estados mayores que la propia cohesión \vertema{3.B.37}.
\item \emph{El \concepto{Fondo de Solidaridad de la Unión Europea}} (2002, tras las inundaciones de Europa central): ayudas tras catástrofes naturales y, desde 2020, emergencias de salud pública. Se financia con la Reserva Europea de Solidaridad, fuera de los techos del marco (hasta 1.016 millones de euros de 2018 al año desde la revisión de 2024)\nocite{UE2024RevisionMFP}.
\item \emph{El Instrumento de Ayuda Preadhesión} (IPA III) prepara a los candidatos, también en desarrollo regional \vertema{3.B.42}.
\item \emph{Instrumentos financieros}: parte de los fondos se da como préstamos, garantías o capital, que se devuelven y se reutilizan, para las inversiones que generan ingresos (las pymes, la eficiencia energética). Crecieron de 1.000 millones en 2000-2006 a 11.500 en 2007-2013\nocite{CE2018Panorama64}.
\item \emph{El \concepto{Banco Europeo de Inversiones}} (BEI) presta para proyectos de interés de la Unión, cofinancia con los fondos y asesora a las regiones. En 2025, el grupo BEI dedicó 42.800 millones, el 48~\% de su financiación en la Unión, a regiones de cohesión\nocite{BEI2026Cohesion} \vertema{3.B.30}.
\item \emph{InvestEU} (Reglamento 2021/523): una garantía del presupuesto de 26.200 millones de euros, gestionada sobre todo por el BEI, para movilizar más de 372.000 millones de inversión. Los Estados pueden aportar a ella parte de sus fondos de cohesión\nocite{UE2021InvestEU}.
\end{itemize}

\begin{anotaciones}
\textbf{Los instrumentos presupuestarios y financieros, de un vistazo}: los fondos de la cohesión comparten un reglamento (el RDC) y se programan a 7 años; los demás instrumentos tienen otros fines, otra gestión o actúan solo ante una emergencia. Importes de la Unión para 2021-2027, en millones de euros corrientes salvo que se diga otra cosa; MCE: Mecanismo «Conectar Europa»; MRR: Mecanismo de Recuperación y Resiliencia; NGEU: NextGenerationEU.
\begin{tablacaja}
\centering
\caption{Instrumentos presupuestarios y financieros de la cohesión, 2021-2027}\label{tab:instrumentos}
\scriptsize
\begin{tabular}{@{}>{\raggedright\arraybackslash}p{2.35cm}>{\raggedright\arraybackslash}p{5.6cm}>{\raggedright\arraybackslash}p{3.15cm}>{\raggedright\arraybackslash}p{2.7cm}@{}}
\toprule
\rowcolor{tceeazulpalido}\textbf{Instrumento} & \textbf{Objetivo} & \textbf{Dotación 2021-2027} & \textbf{Base jurídica}\\
\midrule
\rowcolor{tceeazulpalido}\multicolumn{4}{l}{\textbf{Fondos de la cohesión (gestión compartida; RDC, Reg. 2021/1060)}}\\
\textbf{FEDER} & desequilibrios regionales: innovación, pymes, energía, transporte, ciudades & unos 214.000 & art. 176 TFUE; Reg. 2021/1058\\
\textbf{Fondo de Cohesión} & medio ambiente y redes de transporte en Estados con RNB por habitante $<90$~\% & unos 39.000 (+11.300 al MCE) & art. 177 TFUE; Reg. 2021/1058\\
\textbf{FSE+} & empleo, educación, inclusión y ayuda alimentaria & unos 95.000 & art. 164 TFUE; Reg. 2021/1057\\
\textbf{FTJ} & regiones más afectadas por la transición climática & unos 19.700 (22.000 al inicio) & art. 175 TFUE; Reg. 2021/1056\\
\textbf{Interreg} & cooperación transfronteriza, transnacional e interregional & unos 9.400 (del FEDER) & art. 178 TFUE; Reg. 2021/1059\\
\rowcolor{tceeazulpalido}\multicolumn{4}{l}{\textbf{Fondos de otras políticas con fines parecidos}}\\
\textbf{FEADER} & desarrollo rural (segundo pilar de la PAC) & unos 95.500 (8.100 del NGEU) & art. 43 TFUE; Reg. 2021/2115\\
\textbf{FEMPA} & pesca, acuicultura y zonas costeras & 6.108 & Reg. 2021/1139\\
\textbf{MRR} & reformas e inversiones de los planes nacionales, pagadas por hitos (2021-2026) & hasta 672.500 de 2018 (312.500 en subvenciones) & art. 175 TFUE; Reg. 2021/241\\
\textbf{IPA III} & preparar a los candidatos a la adhesión & 14.162 & art. 212 TFUE; Reg. 2021/1529\\
\rowcolor{tceeazulpalido}\multicolumn{4}{l}{\textbf{Emergencia y solidaridad (fuera de los techos del marco)}}\\
\textbf{Fondo de Solidaridad} & catástrofes naturales y emergencias de salud pública & hasta 1.016 de 2018 al año (Reserva Europea de Solidaridad) & Reg. (CE) 2012/2002\\
\textbf{FEAG} & despedidos en grandes reestructuraciones (más de 200) & hasta 30 de 2018 al año & art. 175 TFUE; Reg. 2021/691\\
\textbf{SURE} & préstamos para la reducción de jornada en la pandemia & 98.400 a 19 Estados (2020-2022) & art. 122 TFUE; Reg. 2020/672\\
\rowcolor{tceeazulpalido}\multicolumn{4}{l}{\textbf{Instrumentos financieros y BEI}}\\
\textbf{Instrumentos financieros} & préstamos, garantías y capital con cargo a los fondos, que se devuelven y se reutilizan & parte de los fondos & arts. 58-62 RDC\\
\textbf{BEI} & préstamos para proyectos de interés de la Unión; asesoramiento a las regiones & 42.800 en regiones de cohesión (2025) & arts. 308-309 TFUE\\
\textbf{InvestEU} & garantía del presupuesto para movilizar inversión & garantía de 26.152 (372.000 movilizados) & Reg. 2021/523\\
\textbf{STEP} & reorientar fondos a tecnologías digitales, limpias y biotecnológicas & 15.200 de la cohesión & Reg. 2024/795\\
\bottomrule
\end{tabular}
\fuente{elaboración propia a partir de los reglamentos citados\nocite{UE2021RDC,UE2021FSE,UE2021FEAG,UE2021InvestEU,UE2021FEMPA,UE2021MRR,UE2021IPA,UE2024RevisionMFP}, la Comisión Europea\nocite{CE2026Cohesiondata,CE2026FEADER,CE2026SURE}, el Parlamento Europeo\nocite{PE2026FichaCohesion} y el BEI\nocite{BEI2026Cohesion}.}
\end{tablacaja}
\end{anotaciones}

% --------------------------------------------------------------------------
\subsection{Los fondos regionales: cómo se gestionan}\label{sec:gestion}

\subsubsection{Gestión compartida, condicionalidades y ejecución}

\begin{itemize}
\item \emph{En gestión compartida}: la Comisión aprueba las estrategias y paga; los Estados y las regiones seleccionan los proyectos, pagan a los beneficiarios y controlan \vertema{3.B.37}. El ciclo:
  \begin{enumerate}
  \item cada Estado negocia con la Comisión un \concepto{acuerdo de asociación}, su estrategia para los 7 años, que debe tener en cuenta las recomendaciones del Semestre Europeo;
  \item los \emph{programas} (nacionales o regionales, por fondo) concretan prioridades, objetivos e indicadores; los aprueba la Comisión;
  \item cada programa tiene una \emph{autoridad de gestión}, una de \emph{auditoría} y un \emph{comité de seguimiento} con los agentes sociales y la sociedad civil (el principio de \emph{asociación} y la \emph{gobernanza multinivel});
  \item la Comisión reembolsa los gastos certificados y vigila; en 2025 se revisan los programas (la revisión intermedia).
  \end{enumerate}
\item \emph{Las condicionalidades}:
  \begin{itemize}
  \item \emph{macroeconómica} (art. 19 RDC): el Consejo puede suspender compromisos o pagos a un Estado que no corrija su déficit excesivo o sus desequilibrios \vertema{3.B.45};
  \item del \emph{Estado de Derecho} (Reglamento 2020/2092) y de la Carta de los Derechos Fundamentales: con ellas se suspendieron fondos a Hungría desde 2022 \vertema{3.B.37}.
  \end{itemize}
\item \emph{La ejecución}: el gran problema de la cohesión es gastar a tiempo. Los programas de 2021-2027 empezaron tarde, porque coincidieron con el MRR, que tenía prioridad por su plazo de 2026:
  \begin{itemize}
  \item en noviembre de 2026, la Unión había pagado el 30~\% de lo programado, y solo el 21~\% sin contar los anticipos; España, el 20~\% y el 12~\%, entre los más bajos\nocite{CE2026Cohesiondata};
  \item los programas de 2014-2020 terminaron pagando el 99,6~\% de lo previsto, pero con años de retraso y concentrando el gasto al final\nocite{CE2026Cohesiondata}.
  \end{itemize}
\end{itemize}

\subsubsection{Retos de futuro}\label{sec:retos-regional}

\begin{itemize}
\item \emph{Los retos de fondo}: gastar a tiempo y bien (la absorción), llegar a las regiones de renta media atrapadas en el estancamiento (\ref{sec:caracterizacion}) y no perder su objetivo, la convergencia, cuando se usa para cada nueva prioridad (\ref{sec:revision-intermedia}).
\item \emph{La propuesta para 2028-2034} (julio de 2025)\nocite{CE2025PropuestaFondo}:
  \begin{itemize}
  \item un solo \emph{Fondo europeo} con la cohesión, la PAC, la pesca, la migración y la seguridad, de 865.000 millones de euros corrientes, repartido en \concepto{planes de asociación nacionales y regionales}: 27 planes y uno de Interreg en vez de unos 540 programas;
  \item al menos 218.000 millones reservados a las regiones menos desarrolladas, con mínimos por Estado, y un mínimo para la PAC;
  \item pagos ligados a hitos y reformas, como en el MRR, y no a costes certificados;
  \item el debate: las regiones (Comité de las Regiones), el Parlamento y los «Amigos de la Cohesión» (con España) temen que el plan nacional deje fuera a las regiones y diluya la cohesión entre otras prioridades; la Comisión defiende la simplificación y el vínculo con las reformas. A noviembre de 2026 el Parlamento aún no ha votado su informe y la negociación del marco sigue abierta \vertema{3.B.37}.
  \end{itemize}
\end{itemize}

\begin{itemize}
\item \emph{Recapitulación}: la política regional se justifica por la tendencia de la integración a concentrar la actividad y se ejecuta con unos 392.000 millones en 2021-2027, en gestión compartida, con reglas de concentración y condiciones. Se ha usado cada vez más para otras prioridades, y su forma después de 2027 está en juego. Pasamos ahora a la otra mitad de la cohesión: las personas.
\end{itemize}

% ==========================================================================
\section{Políticas social y de empleo}\label{sec:social}
\minutos{7}

\begin{itemize}
\item Una vez vista la política que invierte en las regiones, pasamos a la que actúa sobre las personas. Política social y política de empleo se solapan (la mejor política social suele ser el empleo, y las 2 se financian con el FSE+), pero tienen bases jurídicas distintas.
\item \emph{Conclusión anticipada}: la Unión tiene poco dinero y pocas competencias en lo social. Fija mínimos en las condiciones de trabajo, coordina las políticas nacionales con objetivos comunes (el Pilar Europeo de Derechos Sociales y el Semestre) y financia con el FSE+. Su peso ha crecido desde 2017, pero los Estados siguen reservándose los salarios y la protección social, como acaba de recordar el Tribunal de Justicia con la directiva de salarios mínimos.
\end{itemize}

% --------------------------------------------------------------------------
\subsection{Política social}\label{sec:politica-social}

\subsubsection{Justificación}

\begin{itemize}
\item \emph{Su relevancia}: en 2025, 92,7 millones de personas, el 20,9~\% de la población, estaban en riesgo de pobreza o exclusión social, y la Unión dedica a lo social y al empleo unos 95.000 millones del FSE+ en 2021-2027 (\ref{sec:pilar})\nocite{Eurostat2026AROPE,UE2021FSE}.
\item \emph{Las de la política regional}, porque la cohesión es también social.
\item \emph{Evitar el \emph{dumping} social}: en un mercado único, un Estado podría ganar competitividad rebajando sus normas laborales y sociales, y los demás tendrían que seguirle (una «carrera a la baja»). Los mínimos comunes la impiden \vertema{3.B.39}.
\item \emph{La movilidad}: si los trabajadores circulan libremente, hay que coordinar la seguridad social para que no pierdan derechos al cambiar de país.
\item \emph{La legitimidad}: una Unión solo de mercado pierde el apoyo de quienes salen perdiendo con la integración.
\end{itemize}

\subsubsection{Marco jurídico}

\begin{itemize}
\item \emph{Delimitación de competencias}: la política social es \emph{compartida} «en los aspectos definidos» en el TFUE (art. 4.2.b), y la Unión \emph{puede} coordinar las políticas sociales de los Estados (art. 5.3; anexo~\ref{anexo:competencias}). Es una categoría ambigua, fruto de un compromiso entre los Estados que querían una Europa social y los que no\nocite{UE2016TFUE}. La \emph{cláusula social horizontal} (art. 9) obliga a tener en cuenta el empleo, la protección social y la lucha contra la exclusión en todas las políticas.
\item \emph{Marco jurídico} (título X del TFUE, arts. 151-161):
  \begin{itemize}
  \item \emph{los objetivos} (art. 151): el fomento del empleo, la mejora de las condiciones de vida y de trabajo, una protección social adecuada, el diálogo social, el desarrollo de los recursos humanos y la lucha contra las exclusiones, con referencia a la Carta Social Europea (1961) y a la Carta comunitaria de 1989;
  \item \emph{lo que puede hacer la Unión} (art. 153): directivas de mínimos en salud y seguridad en el trabajo, condiciones de trabajo, información y consulta de los trabajadores, integración de los excluidos del mercado laboral, igualdad entre hombres y mujeres, seguridad social y protección social de los trabajadores, protección en caso de despido, representación colectiva y condiciones de empleo de los nacionales de terceros países; y medidas de cooperación, sin armonización, contra la exclusión y para modernizar la protección social;
  \item \emph{lo que no puede hacer}: las normas sociales no se aplican a \emph{las remuneraciones, el derecho de asociación y sindicación, el derecho de huelga ni el de cierre patronal} (art. 153.5), y no pueden afectar a los principios fundamentales ni al equilibrio financiero de los sistemas nacionales de seguridad social (art. 153.4).
  \end{itemize}
\item \emph{Toma de decisiones}:
  \begin{itemize}
  \item las directivas de mínimos se aprueban por el \emph{procedimiento legislativo ordinario} (mayoría cualificada en el Consejo), previa consulta al Comité Económico y Social y al de las Regiones;
  \item por \emph{unanimidad} del Consejo, previa consulta al Parlamento, las de seguridad social y protección social de los trabajadores, protección en caso de despido, representación colectiva y nacionales de terceros países. Salvo para la seguridad social, el Consejo puede pasar por unanimidad al procedimiento ordinario (pasarela);
  \item \emph{los interlocutores sociales} (arts. 154-155): la Comisión los consulta antes de proponer, y sus acuerdos pueden convertirse en directivas (así nacieron las de permisos parentales, trabajo a tiempo parcial y contratos de duración determinada). El \concepto{Comité de Protección Social} (art. 160) es el órgano consultivo del Consejo.
  \end{itemize}
\end{itemize}

\subsubsection{Antecedentes}

\begin{itemize}
\item \emph{1957}: el Tratado de Roma solo prevé la igualdad de retribución entre hombres y mujeres (hoy art. 157), la coordinación de la seguridad social de los trabajadores migrantes y el \concepto{Fondo Social Europeo} (FSE), para la reconversión profesional; es el fondo más antiguo.
\item \emph{1961}: Carta Social Europea, del Consejo de Europa.
\item \emph{1986-1989}: el Acta Única permite directivas de salud y seguridad por mayoría cualificada; en 1989, la Carta comunitaria de los derechos sociales fundamentales de los trabajadores, sin el Reino Unido y sin valor jurídico.
\item \emph{1992}: el Protocolo social de Maastricht, del que se excluye el Reino Unido, amplía la mayoría cualificada y el diálogo social.
\item \emph{1997}: Ámsterdam integra el Protocolo en el Tratado (el Reino Unido se suma) y crea el título de empleo.
\item \emph{2000-2009}: la Carta de los Derechos Fundamentales de la Unión (Niza, 2000), con un capítulo de solidaridad; desde Lisboa (2009), con el mismo valor jurídico que los Tratados.
\item \emph{2010}: la Estrategia Europa 2020, con objetivos sociales: tasa de empleo del 75~\%, 20 millones de personas menos en riesgo de pobreza, abandono escolar por debajo del 10~\% y un 40~\% de jóvenes con estudios superiores.
\end{itemize}

\subsubsection{Situación actual: el Pilar Europeo de Derechos Sociales}\label{sec:pilar}

\begin{itemize}
\item \emph{Qué es}: proclamado por el Parlamento, el Consejo y la Comisión en la cumbre de Gotemburgo (17 de noviembre de 2017), recoge \emph{20 principios} en 3 capítulos: igualdad de oportunidades y acceso al mercado de trabajo; condiciones de trabajo justas; protección e inclusión social. No es vinculante: es una brújula para la legislación de la Unión y para las políticas nacionales.
\item \emph{Su plan de acción} (marzo de 2021), refrendado en la Cumbre Social de Oporto (mayo de 2021), fija 3 objetivos para 2030\nocite{CE2021PlanPilar}:
  \begin{itemize}
  \item una \emph{tasa de empleo} de al menos el 78~\% entre los 20 y los 64 años;
  \item que al menos el 60~\% de los adultos se \emph{forme} cada año;
  \item al menos \emph{15 millones de personas menos en riesgo de pobreza o exclusión social} que en 2019, 5 millones de ellos niños\nocite{CE2021PlanPilar}.
  \end{itemize}
\item \emph{El indicador de pobreza} es el \concepto{AROPE}, siglas del inglés \emph{at risk of poverty or social exclusion} (en riesgo de pobreza o exclusión social). Eurostat lo define como el porcentaje de la población que está en al menos una de 3 situaciones\nocite{EurostatGlosarioAROPE}:
  \begin{itemize}
  \item \emph{riesgo de pobreza}: renta disponible equivalente, después de las transferencias sociales, por debajo del 60~\% de la mediana nacional;
  \item \emph{privación material y social severa}: no poder permitirse al menos 7 de 13 conceptos básicos (pagar a tiempo las facturas, una semana de vacaciones, una comida con carne o pescado (o su equivalente vegetariano) cada 2 días, mantener la casa caliente, afrontar un gasto imprevisto, tener coche, sustituir los muebles estropeados, entre otros);
  \item \emph{muy baja intensidad de empleo}: vivir (con menos de 65 años) en un hogar cuyos miembros de 18 a 64 años trabajaron el 20~\% o menos de su tiempo potencial el año anterior.
  \end{itemize}
\item \emph{Cómo van} (gráfico~\ref{fig:pilar}):
  \begin{itemize}
  \item el empleo, bien: el 76,1~\% en 2025, con 16 Estados ya por encima del 78~\%; España, en el 72,4~\%, con un objetivo nacional del 76~\%\nocite{Eurostat2026Empleo,Consejo2026EMCO};
  \item la pobreza, mal: 92,7 millones de personas en riesgo en 2025 (el 20,9~\%), apenas 2,8 millones menos que en 2019, y los niños en riesgo no bajan. España está en el 25,7~\%, de los más altos de la Unión, con un 34~\% en los niños\nocite{Eurostat2026AROPE};
  \item la formación de adultos, a 13 puntos del objetivo.
  \end{itemize}
\item \emph{La renovación del Pilar} (julio de 2026): la Comisión mantiene los objetivos de 2030 y pone 3 prioridades, la asequibilidad (vivienda, cuidados, una futura ley de empleos de calidad), la reducción de las desigualdades y la inteligencia artificial en el trabajo. La acompañan la primera \emph{Estrategia contra la Pobreza} (mayo de 2026, con la meta de erradicarla en 2050) y el \emph{Plan Europeo de Vivienda Asequible} (diciembre de 2025)\nocite{CE2026Pilar,CE2026Pobreza}.
\end{itemize}

\begin{figure}[H]
\centering
\caption{Los objetivos sociales para 2030: empleo y pobreza}\label{fig:pilar}
\grafico[0.9\textwidth]{pilar-social}
\fuente{elaboración propia a partir de Eurostat\nocite{Eurostat2026Empleo,Eurostat2026AROPE} (tasa de empleo de 20 a 64 años; personas en riesgo de pobreza o exclusión social, serie de seguimiento del objetivo).}
\end{figure}

\subsubsection{Instrumentos}

\begin{itemize}
\item \emph{Normativos}: las directivas de mínimos. Las más recientes:
  \begin{itemize}
  \item la \concepto{Directiva de salarios mínimos adecuados} (2022/2041): no fija un salario mínimo europeo, sino un procedimiento. Los Estados con salario mínimo legal deben actualizarlo al menos cada 2 años con criterios claros, con valores de referencia como el 60~\% de la mediana o el 50~\% de la media de los salarios brutos; los que tengan una cobertura de la negociación colectiva inferior al 80~\% deben aprobar un plan para aumentarla\nocite{UE2022SalarioMinimo}. El \emph{Tribunal de Justicia} (Gran Sala, 11 de noviembre de 2025, asunto C-19/23, Dinamarca contra Parlamento y Consejo) la ha mantenido en lo esencial, pero ha anulado la lista obligatoria de criterios para fijar el salario (art. 5.2) y la regla que impedía que la indexación automática lo bajara, porque son una injerencia directa en la fijación de las remuneraciones, excluida por el art. 153.5\nocite{TJUE2025SalarioMinimo};
  \item la \emph{Directiva de transparencia salarial} (2023/970), para cerrar la brecha de género, que debía transponerse antes del 7 de junio de 2026;
  \item la \emph{Directiva de trabajo en plataformas digitales} (2024/2831): presunción de laboralidad cuando hay control y dirección, con la carga de la prueba en la plataforma, y límites a la gestión algorítmica (supervisión humana, transparencia). Se transpone antes del 2 de diciembre de 2026; España se adelantó con la «ley rider» (Ley 12/2021)\nocite{UE2024Plataformas}.
  \end{itemize}
\item \emph{De coordinación}: el \emph{método abierto de coordinación} (objetivos comunes, indicadores y aprendizaje mutuo) y el \emph{cuadro de indicadores sociales} del Semestre (\ref{sec:empleo}); la \emph{Garantía Infantil Europea} (Recomendación de 2021), para que los niños en riesgo tengan acceso a educación, sanidad, nutrición y vivienda. Su indicador clave es el AROPE (\ref{sec:pilar}).
\item \emph{Financieros}: el \concepto{Fondo Social Europeo Plus} (FSE+), de unos 99.000 millones de euros corrientes en 2021-2027 (87.995 millones de 2018; unos 95.000 en gestión compartida), que integra el antiguo FSE, la Iniciativa de Empleo Juvenil, el \emph{Fondo de Ayuda Europea para las Personas Más Desfavorecidas} (FEAD) y el programa de empleo e innovación social (EaSI)\nocite{UE2021FSE}:
  \begin{itemize}
  \item al menos el 25~\% para la inclusión social y el 3~\% para la privación material (lo que antes hacía el FEAD: alimentos y bienes básicos);
  \item el 12~\% para el empleo juvenil en los Estados con muchos jóvenes que ni estudian ni trabajan, y el 5~\% para la pobreza infantil en los que superan la media, como España.
  \end{itemize}
\end{itemize}

% --------------------------------------------------------------------------
\subsection{Política de empleo}\label{sec:empleo}

\subsubsection{Justificación}

\begin{itemize}
\item \emph{Eficiencia}: mejores servicios de empleo y más formación casan mejor la oferta con la demanda de trabajo y reducen el paro estructural \vertema{3.A.26}.
\item \emph{Protección}: la política de empleo amortigua las pérdidas de quienes sufren el ajuste (y aquí se toca con la social).
\item \emph{Coordinación en la unión monetaria}: sin tipo de cambio propio, el ajuste a una perturbación asimétrica recae en los salarios y el empleo; mercados de trabajo flexibles y una protección que no frene la movilidad hacen que el euro funcione mejor \vertema{3.B.16}.
\end{itemize}

\subsubsection{Marco jurídico}

\begin{itemize}
\item \emph{Delimitación de competencias}: solo de \emph{coordinación}: la Unión toma medidas para coordinar las políticas de empleo, en particular con orientaciones (art. 5.2 TFUE; anexo~\ref{anexo:competencias}). El empleo es una «cuestión de interés común» (art. 146).
\item \emph{Marco jurídico} (título IX del TFUE, arts. 145-150; su objetivo, en el art. 145): una estrategia coordinada para una mano de obra cualificada, formada y adaptable y unos mercados de trabajo que respondan a los cambios económicos.
\item \emph{Toma de decisiones} (art. 148): cada año el Consejo Europeo examina la situación con el \emph{informe conjunto sobre el empleo}; el Consejo aprueba \emph{orientaciones} para las políticas de empleo, examina los informes nacionales y puede dirigir \emph{recomendaciones} a cada Estado. Las medidas de fomento (art. 149) se aprueban por el procedimiento ordinario y excluyen toda armonización. El \concepto{Comité de Empleo} (art. 150) asesora al Consejo.
\end{itemize}

\subsubsection{Evolución histórica y situación actual}

\begin{itemize}
\item \emph{La Estrategia Europea de Empleo} (Cumbre de Luxemburgo, noviembre de 1997): primer ejemplo del método abierto de coordinación, con 4 pilares (empleabilidad, espíritu de empresa, adaptabilidad e igualdad de oportunidades).
\item \emph{La Estrategia de Lisboa} (2000) y \emph{Europa 2020} (2010) le pusieron objetivos cuantitativos; la Comisión promovió la \concepto{flexiseguridad}: mercados flexibles para contratar y reorganizar, con seguridad para el trabajador entendida como empleabilidad y prestaciones, no como un puesto fijo (el modelo danés).
\item \emph{Desde 2011}, la estrategia está dentro del \concepto{Semestre Europeo} \vertema{3.B.45}:
  \begin{itemize}
  \item en otoño, la Comisión propone el \emph{informe conjunto sobre el empleo}, con el cuadro de indicadores sociales y, desde 2024, el \emph{marco de convergencia social}, que analiza a fondo a los Estados con más riesgos (en 2026, 9 Estados, España entre ellos);
  \item las \emph{orientaciones de empleo} (4, la última vez mantenidas por la Decisión 2025/2254) guían los programas nacionales de reformas;
  \item en primavera, las \emph{recomendaciones específicas por país}, que también guían la programación de los fondos (\ref{sec:gestion}).
  \end{itemize}
\item \emph{La Garantía Juvenil} (2013, reforzada en 2020): todo joven de hasta 29 años debe recibir en 4 meses una oferta de empleo, formación o prácticas. El porcentaje de jóvenes que ni estudian ni trabajan en la Unión ha bajado al 11~\% en 2025, con el objetivo del 9~\% en 2030\nocite{Consejo2026EMCO}.
\end{itemize}

\subsubsection{Instrumentos}

\begin{itemize}
\item \emph{El FSE+}, que financia sobre todo políticas activas: formación, orientación, ayudas a la contratación y movilidad.
\item \emph{El \concepto{Fondo Europeo de Adaptación a la Globalización}} (FEAG, 2007; reformado en 2021): ayudas puntuales a los trabajadores despedidos en reestructuraciones grandes (más de 200 despidos), por la globalización, la transición digital o la verde. Se moviliza fuera de los techos del marco, con un máximo de 30 millones de euros al año (a precios de 2018) y una cofinanciación del 60~\% al 85~\%\nocite{UE2021FEAG}. A diferencia del FSE+, que programa a 7 años, el FEAG reacciona a casos concretos.
\item \emph{EURES}, la red de servicios de empleo para la movilidad entre Estados.
\item \emph{El SURE} (2020-2022): 98.400 millones de euros en préstamos a 19 Estados, financiados con deuda de la Unión, para pagar los regímenes de reducción de jornada (en España, los ERTE, con 21.324 millones). En 2020 sostuvo a unos 31,5 millones de trabajadores y 2,5 millones de empresas\nocite{CE2026SURE}. Su valoración como estabilizador está en otro lugar \vertema{3.B.45} \vertema{4.A.16}.
\end{itemize}

\subsubsection{Retos de futuro}\label{sec:retos-social}

\begin{itemize}
\item \emph{La pobreza}: es el objetivo de 2030 que peor va, sobre todo la infantil; la Estrategia contra la Pobreza de 2026 tendrá que demostrar que la coordinación basta\nocite{CE2026Pobreza}.
\item \emph{Los límites de la competencia}: la sentencia sobre los salarios mínimos marca hasta dónde puede llegar la Unión en las remuneraciones, y cualquier paso más (un seguro de desempleo europeo, un salario mínimo común) exigiría reformar los Tratados o usar otras bases, como el art. 122 en el SURE.
\item \emph{El trabajo que cambia}: el envejecimiento y la falta de mano de obra cualificada, la transición verde y digital, las plataformas y la inteligencia artificial, que el Pilar renovado pone entre sus prioridades\nocite{CE2026Pilar}.
\item \emph{La financiación}: la propuesta para 2028-2034 integra el FSE+ en los planes nacionales y regionales; el Parlamento y los interlocutores sociales temen que el gasto social se diluya entre otras prioridades\nocite{CE2025PropuestaFondo}.
\end{itemize}

\begin{itemize}
\item \emph{Recapitulación}: en lo social y el empleo, la Unión pone mínimos, objetivos y dinero, pero no los salarios ni la protección social. Queda por ver si todo esto, junto con la política regional, ha acercado las rentas.
\end{itemize}

% ==========================================================================
\section{El proceso de convergencia real en la Unión Europea}\label{sec:convergencia}
\minutos{8,5}

\begin{itemize}
\item Vistas las políticas, queda la prueba de los resultados. Responderemos a 2 preguntas: ¿se han acercado las rentas de los países y las regiones de la Unión?, y ¿qué parte se debe a los fondos? Nos centramos en la renta por habitante en PPA, aunque la convergencia también se mide en productividad, empleo o paro.
\item \emph{Conclusión anticipada}: ha habido convergencia, rápida, entre países y entre regiones, pero casi toda la explica el Este; el Sur, España incluida, se ha alejado desde 2008, y muchas regiones de renta media están estancadas. Los fondos ayudan a crecer, en torno al umbral del 75~\%, pero su efecto depende de la capacidad de cada región para aprovecharlos.
\end{itemize}

% --------------------------------------------------------------------------
\subsection{Caracterización teórica de la convergencia}\label{sec:caracterizacion}

\paragraph{a) Idea}

\begin{itemize}
\item La teoría del crecimiento predice que una economía crece más deprisa cuanto más lejos está de su estado estacionario, porque el capital rinde más donde es escaso. Si las regiones de la Unión comparten tecnología, instituciones y mercado, sus estados estacionarios se parecerán y las pobres deberían alcanzar a las ricas \vertema{3.A.43}.\footnote{\autor{Robert Solow} recibió el Premio Nobel de Economía en 1987 «por sus contribuciones a la teoría del crecimiento económico»\nocite{Nobel1987Solow}.}
\item Para comprobarlo se usan las nociones de convergencia $\beta$ y $\sigma$, de \autor{Xavier Sala-i-Martín} (tesis de 1990) y \textcite{BarroSalaiMartin1991}, que aquí resumimos y aplicamos a las regiones de la Unión \vertema{3.A.45}.
\end{itemize}

\paragraph{b) Modelo}

\subparagraph{i) Supuestos}

\begin{enumerate}
\item Cada región $i$ se comporta como una economía de \autor{Solow}: rendimientos decrecientes del capital y progreso técnico exógeno común.
\item Su crecimiento depende de la distancia a su estado estacionario, que depende de unos determinantes $\vect{z}$ (tasa de inversión, crecimiento de la población, capital humano).
\item Las perturbaciones regionales son transitorias.
\end{enumerate}

\subparagraph{ii) Desarrollo}

\begin{itemize}
\item \emph{Paso 1: la convergencia $\beta$}. Hay \concepto{convergencia $\beta$} si las regiones con menos renta por habitante crecen más que las de más renta. Se contrasta con una regresión entre regiones del crecimiento sobre la renta inicial:
  \funcionclave[eq:convergencia-beta]{\gamma^{y}_{t,t-1}=\beta_0+\beta_1\ln y_{t-1}+\vect{\eta}\cdot\vect{z}_{t-1}+u_t}
  \begin{itemize}
  \item es \emph{absoluta} si $\beta_1<0$ sin controles ($\vect{\eta}=\vect{0}$), y \emph{condicional} si $\beta_1<0$ al controlar por $\vect{z}_{t-1}$;
  \item la \emph{velocidad de convergencia} $\lambda$ es la fracción de la distancia al estado estacionario que se cierra cada año: con un periodo de $T$ años, $\beta_1=-(1-e^{-\lambda T})/T$. En el modelo de \autor{Solow}, $\lambda=(1-\alpha)(n+g+\delta)$.
  \end{itemize}
\item \emph{Paso 2: la convergencia $\sigma$}. Hay \concepto{convergencia $\sigma$} si la dispersión de las rentas por habitante (la varianza de su logaritmo) se reduce con el tiempo:
  \begin{equation}
  \sigma^{2}_{y,t-1}>\sigma^{2}_{y,t}.\label{eq:convergencia-sigma}
  \end{equation}
\item \emph{Paso 3: la relación entre las 2}. La $\beta$ es necesaria pero no suficiente para la $\sigma$: una región pobre puede crecer tanto que adelante a la rica y deje más distancia que al principio, y las perturbaciones nuevas pueden volver a separar a las regiones aunque las pobres crezcan más. Los 3 casos con 2 economías se dibujan aparte \vertema{3.A.45}.
\end{itemize}

\subparagraph{iii) Implicaciones}

\begin{itemize}
\item Si hay convergencia absoluta, el mercado único bastaría y la política regional solo la aceleraría.
\item Si solo hay convergencia condicional, cada región va a su propio estado estacionario, y una región con poco capital humano o malas instituciones se queda atrás \emph{para siempre}: la política debe cambiar los determinantes $\vect{z}$, no solo dar renta.
\item Con $\lambda\simeq2$~\%, la mitad de la distancia tarda unos 35 años en cerrarse: la convergencia es cosa de generaciones.
\end{itemize}

\paragraph{c) Extensiones}

\begin{itemize}
\item \emph{Clubes de convergencia}: las regiones pueden converger entre grupos distintos, no hacia una media común. \textcite{Quah1996} encuentra en Europa una distribución con 2 picos, y que lo que pesa es el país y los vecinos, no la Unión en su conjunto.
\item \emph{La trampa del desarrollo regional}: \textcite{DiemerEtAl2022} llaman así a las regiones que crecen menos que la Unión, su país y su propio pasado. La Comisión identifica en esa situación al 28~\% de la población de la Unión (2021), sobre todo en regiones de renta media y alta, y añade una \emph{trampa del talento}: 46 regiones con más de 70 millones de habitantes pierden población joven y cualificada\nocite{CE2024Informe9}.
\end{itemize}

\paragraph{d) Evidencia empírica}

\begin{itemize}
\item \emph{La clásica}: \textcite{BarroSalaiMartin1991} y \textcite{SalaiMartin1996Regional} encuentran entre regiones europeas una velocidad del 2~\% anual, igual que entre los estados de Estados Unidos y las prefecturas de Japón.
\item \emph{Una estimación propia con los datos de hoy}: con las 232 regiones NUTS 2 de la Unión y su PIB por habitante en PPA relativo a la media, entre 2000 y 2024 (gráfico~\ref{fig:convergencia-regional})\nocite{Eurostat2026Regional}:
  \begin{itemize}
  \item hay convergencia $\beta$ absoluta y fuerte: $\hat\beta_1=-1{,}99$ (error típico 0,12), con un $R^2$ de 0,55. Una región con la mitad de renta que otra crece 1,4 puntos más al año ($1{,}99\times\ln2$), y la velocidad implícita es $\lambda\simeq2{,}7$~\% anual, algo mayor que la clásica;
  \item pero la explican las regiones del Este, que parten de muy abajo y crecen 2-4 puntos más que la media. Entre las regiones del resto de la Unión la pendiente es nula ($+0{,}02$, con un $R^2$ de 0): las pobres del Sur no crecen más que las ricas del Norte, y muchas crecen menos;
  \item las regiones españolas están casi todas por debajo de cero: Canarias pasó del 91~\% de la media en 2000 al 72~\% en 2024, la Comunidad Valenciana del 93~\% al 77~\% y Andalucía del 72~\% al 69~\%; solo Extremadura y Galicia mejoran algo.
  \end{itemize}
\item \emph{La convergencia $\sigma$}: la dispersión baja entre Estados (de 0,56 en 2000 a 0,31 en 2024) y entre regiones (de 0,46 a 0,34), pero entre regiones se para de 2008 a 2016 (gráfico~\ref{fig:sigma}).
\end{itemize}

\begin{figure}[H]
\centering
\caption{Convergencia $\beta$ entre las regiones de la Unión, 2000-2024}\label{fig:convergencia-regional}
\grafico[0.85\textwidth]{convergencia-regional}
\fuente{elaboración propia a partir de Eurostat\nocite{Eurostat2026Regional} (PIB por habitante en PPA de las regiones NUTS 2, UE-27 = 100). Regresión por mínimos cuadrados de $\frac{100}{24}\ln(y_{2024}/y_{2000})$ sobre $\ln y_{2000}$; línea discontinua, la misma regresión solo con las regiones del resto de la Unión. Regiones del Este: Estados que entraron entre 2004 y 2013, salvo Chipre y Malta.}
\end{figure}

\begin{figure}[H]
\centering
\caption{Convergencia $\sigma$ en la Unión: entre Estados y entre regiones, 2000-2024}\label{fig:sigma}
\grafico[0.7\textwidth]{sigma-regional}
\fuente{elaboración propia a partir de Eurostat\nocite{Eurostat2026Regional,Eurostat2026PPA} (desviación típica del logaritmo del PIB por habitante en PPA; 224 regiones NUTS 2 con datos todos los años y 27 Estados, sin ponderar).}
\end{figure}

\paragraph{e) Valoración}

\begin{itemize}
\item El modelo da una forma sencilla y comparable de medir la convergencia, y los datos de la Unión lo confirman a medias: convergen las regiones que estaban muy lejos, y se estancan las de renta media.
\item Tiene límites conocidos: la regresión a la media de \autor{Galton} (una $\beta$ negativa puede salir solo de perturbaciones transitorias), el sesgo de las variables omitidas y que el PPA sirve para comparar niveles en un año, no tanto para medir crecimientos. Por eso conviene mirar también la evidencia descriptiva y, para los fondos, métodos que identifiquen la causa.
\end{itemize}

% --------------------------------------------------------------------------
\subsection{Evidencia sobre la convergencia en la Unión Europea}\label{sec:evidencia}

\begin{itemize}
\item \emph{Hasta 2008, convergencia}: la ampliación al Este coincide con un crecimiento muy rápido de los nuevos Estados. En las regiones, la renta media del 20~\% más rico era unas 3,5 veces la del 20~\% más pobre en 2000, y 2,8 veces en 2008\nocite{CE2022Informe8}.
\item \emph{2008-2013, la crisis la frena}: las diferencias de paro se disparan (en 2013, 27 regiones tenían más del 20~\% de paro, frente a 5 en 2008) y el Sur retrocede\nocite{CE2022Informe8}.
\item \emph{Desde 2014, convergencia a 2 velocidades} (gráfico~\ref{fig:pib-pc-paises})\nocite{Eurostat2026PPA}:
  \begin{itemize}
  \item \emph{el Este} sigue acercándose: los nuevos Estados pasan del 45~\% de la media en 1995 a cerca del 80~\% hoy\nocite{CE2024Informe9}. Chequia ya iguala a España (92~\%), Polonia (81~\%) ha alcanzado a Portugal y Lituania la supera;
  \item \emph{el Sur} se aleja: Italia baja del 122~\% al 96~\% desde 2000, Grecia está en el 68~\% y Portugal no se mueve del 81~\%. Un tercio de las regiones de la Unión no había recuperado en 2023 su renta por habitante de 2008, y entre ellas están toda Grecia, el 80~\% de Italia y un tercio de España\nocite{CE2024Informe9}.
  \end{itemize}
\item \emph{España}: del 91~\% de la media en 1995 al 105~\% en 2006, con el auge del crédito y la construcción; cae al 90~\% tras la crisis, al 82~\% en 2020, y vuelve al 92~\% en 2025. Su convergencia anterior a 2008 no era sostenible: venía del empleo en sectores de baja productividad, no de la productividad, que apenas creció \vertema{4.A.28} \vertema{4.A.17}.
\item \emph{El patrón}: convergencia entre países y, dentro de cada país, divergencia a favor de las capitales y las regiones con servicios avanzados; el Este converge por la integración en las cadenas de valor alemanas y la inversión extranjera, y el Sur sufre la competencia de Asia en manufacturas, la crisis del euro y la baja productividad\nocite{IammarinoEtAl2019}.
\end{itemize}

\begin{figure}[H]
\centering
\caption{España frente al Sur y el Este: PIB por habitante en PPA, 1995-2025 (UE-27 = 100)}\label{fig:pib-pc-paises}
\grafico[0.97\textwidth]{pib-pc-paises}
\fuente{elaboración propia a partir de Eurostat\nocite{Eurostat2026PPA}.}
\end{figure}

\begin{anotaciones}
\textbf{El mapa de la renta regional}: en 2024, las regiones por debajo del 75~\% de la media se concentran en Bulgaria, Rumanía, Hungría, Grecia, el sur de Italia, el sur de España y las ultraperiféricas, y las de más del 150~\%, en las capitales (Luxemburgo, Dublín, Bruselas, Praga, París, Copenhague, Ámsterdam, Estocolmo, Budapest, Bucarest, Varsovia y Bratislava) y en polos como Múnich, Stuttgart y Hamburgo. Entre 2000 y 2024 ganaron renta relativa casi todas las regiones del Este y de Irlanda; la perdieron casi todas las de Italia, Francia, Grecia, Suecia y Finlandia, y la mayoría de las de España.
\begin{figuracaja}
\caption{El PIB por habitante de las regiones de la Unión: nivel en 2024 y variación desde 2000 (PPA, UE-27 = 100)}\label{fig:mapa-pib}
\grafico[\linewidth]{mapa-pib-regional}
\fuente{elaboración propia a partir de Eurostat\nocite{Eurostat2026Regional} (regiones NUTS 2 de 2024; entre paréntesis, el número de regiones). Contornos: GISCO, Eurostat (© EuroGeographics para los límites administrativos)\nocite{GISCO2026NUTS}.}
\end{figuracaja}
\end{anotaciones}

% --------------------------------------------------------------------------
\subsection{Evidencia sobre el papel de la política regional}\label{sec:papel}

\begin{itemize}
\item La convergencia del Este coincide con los fondos, pero también con la entrada en el mercado único y la llegada de capital extranjero. ¿Cuánto se debe a los fondos? Hay 2 maneras de responder: los modelos macroeconómicos, que simulan la economía con y sin fondos, y la evaluación de impacto, que compara regiones parecidas que los reciben y que no.
\end{itemize}

\subsubsection{El efecto de los fondos: la regresión en discontinuidad}\label{sec:rdd}

\paragraph{a) Idea}

\begin{itemize}
\item Comparar el crecimiento de las regiones con y sin fondos no sirve: las que los reciben son, por definición, más pobres, y crecerían distinto de todos modos (el sesgo de selección).
\item Pero la regla de reparto da un experimento natural: una región con el 74~\% de la renta media recibe muchos más fondos que una con el 76~\%, y por lo demás son casi iguales. \textcite{BeckerEggerVonEhrlich2010} usan ese salto en el umbral del 75~\% para medir el efecto de los fondos del antiguo objetivo 1 con una \concepto{regresión en discontinuidad} (\emph{regression discontinuity design}).
\end{itemize}

\paragraph{b) Modelo}

\subparagraph{i) Supuestos}

\begin{enumerate}
\item La variable de asignación $x_i$ es la renta por habitante de la región en \% de la media en los años de referencia, y el tratamiento es $D_i=\mathbb{1}[x_i<75]$.
\item Las regiones no pueden manipular $x_i$ para quedar justo por debajo del umbral.
\item Sin fondos, el crecimiento $g_i$ variaría de forma continua con $x_i$: todo lo demás que influye en el crecimiento no salta en el 75~\%.
\end{enumerate}

\subparagraph{ii) Desarrollo}

\begin{itemize}
\item \emph{Paso 1}: se estima el crecimiento posterior en función de la distancia al umbral, a los 2 lados, con un salto en el umbral:
  \funcionclave[eq:rdd]{g_i=\alpha+\tau D_i+f(x_i-75)+\varepsilon_i}
  donde $f(\cdot)$ es una función flexible (un polinomio, que puede ser distinto a cada lado).
\item \emph{Paso 2}: $\tau$ es el salto del crecimiento en el umbral (gráfico~\ref{fig:rdd}): como las regiones a un lado y al otro son comparables, es el \emph{efecto causal} de los fondos en las regiones cercanas al 75~\%.
\item \emph{Paso 3}: como no todas las regiones por debajo reciben el tratamiento ni todas las de encima se quedan sin él (régimen transitorio, excepciones), se usa una versión \emph{borrosa} (\emph{fuzzy}), que divide el salto del crecimiento por el salto de la probabilidad de tratamiento.
\end{itemize}

\begin{figure}[H]
\centering
\caption{La regresión en discontinuidad en el umbral del 75~\%}\label{fig:rdd}
\grafico[0.68\textwidth]{rdd-umbral}
\fuente{elaboración propia a partir de \textcite{BeckerEggerVonEhrlich2010}. Ilustrativo: los puntos no son datos. $x_i$: PIB por habitante en PPA de la región en \% de la media de la Unión; $g_i$: crecimiento posterior.}
\end{figure}

\subparagraph{iii) Implicaciones}

\begin{itemize}
\item \textcite{BeckerEggerVonEhrlich2010} encuentran que las regiones del objetivo 1 crecen más en PIB por habitante (en torno a 1,6 puntos más de crecimiento anual medio durante el periodo de programación), pero \emph{no en empleo}; con un cálculo sencillo de costes y beneficios, concluyen que los fondos no solo son eficaces, sino también eficientes.
\item Es decir, el pan sirve: los fondos sí suben la renta de las regiones que los reciben.
\end{itemize}

\paragraph{c) Extensiones}

\begin{itemize}
\item \emph{La capacidad de absorción}: \textcite{BeckerEggerVonEhrlich2013} muestran que solo en torno al 30~\% de las regiones convierte los fondos en más crecimiento: las que tienen suficiente capital humano y buenas instituciones. Los fondos rinden más donde la región tiene con qué aprovecharlos.
\item \emph{La persistencia}: \textcite{BeckerEggerVonEhrlich2018} encuentran que el efecto se deshace en buena parte cuando la región pierde la condición de objetivo 1, y que fue más débil durante la crisis.
\item \emph{Otras estimaciones}: \textcite{PellegriniEtAl2013} confirman con el mismo método un efecto positivo en el crecimiento en 1994-2006.
\end{itemize}

\paragraph{d) Evidencia empírica}

\begin{itemize}
\item \emph{Los modelos macroeconómicos de la Comisión} (QUEST, RHOMOLO; antes HERMIN y EcoMod) estiman que cada euro invertido en 2014-2027 genera 1,3 euros de PIB en 2030 y casi el triple a largo plazo, con efectos mayores en los grandes receptores\nocite{CE2024Informe9}. En 2000-2006, HERMIN estimaba que el PIB de los nuevos Estados habría sido entre un 5~\% y un 10~\% menor sin la cohesión, la mitad por efectos de demanda y la otra mitad por oferta.
\item \emph{En España}: \textcite{delaFuenteVives1995} encuentran que el FEDER, repartido con criterio redistributivo, tuvo un efecto significativo aunque limitado en la convergencia regional, mientras que la inversión pública nacional en infraestructuras apenas contribuyó; trabajos posteriores atribuyen a los fondos en torno a 1/5 o 1/3 de la convergencia de las regiones españolas en los 90 \vertema{4.A.17}.
\item \emph{Los escépticos}: \textcite{BoldrinCanova2001} no encuentran que los fondos reduzcan las disparidades y los ven como una política sobre todo redistributiva, ligada a la negociación política.
\end{itemize}

\paragraph{e) Valoración}

\begin{itemize}
\item La regresión en discontinuidad es la evidencia más creíble: compara regiones casi iguales y no depende de un modelo de toda la economía.
\item Pero mide un efecto \emph{local}, el de las regiones en torno al 75~\%, que no tiene por qué valer para las muy pobres ni para las de transición; y mide el crecimiento en los años de los fondos, no si la región se ha transformado.
\item Junto con la capacidad de absorción, apunta a lo que falla en el Sur: el dinero llega, pero sin capital humano, sin instituciones que seleccionen bien los proyectos y sin estrategias de especialización, se queda en gasto.
\end{itemize}

% --------------------------------------------------------------------------
\subsection{Dificultades de medición}\label{sec:medicion}

\begin{itemize}
\item \emph{Muchos factores a la vez}: la convergencia depende del mercado único, del euro, de la inversión extranjera, de las políticas nacionales y del ciclo; separar el efecto de los fondos es difícil (variables omitidas).
\item \emph{El contrafactual}: hay que saber qué habría pasado sin los fondos y qué coste de oportunidad tienen:
  \begin{itemize}
  \item \emph{expulsión} (\emph{crowding out}): los fondos pueden sustituir a la inversión nacional que se habría hecho igual (por eso existe la adicionalidad) o a la privada, y se financian con impuestos que distorsionan;
  \item \emph{atracción} (\emph{crowding in}): si el proyecto complementa al capital privado (una carretera, una red de banda ancha), sube su rentabilidad y atrae más inversión.
  \end{itemize}
\item \emph{Los efectos fuera de la región}: parte del gasto vuelve a los contribuyentes netos, en importaciones y en contratos para sus empresas, y la mejora de la periferia beneficia también al centro \vertema{3.B.37}.
\item \emph{La medida}: la renta por habitante en PPA es buena para comparar niveles, pero en las regiones con mucha conmutación (Bruselas, Luxemburgo, Praga) el PIB se produce donde no se vive; por eso se miran también la renta disponible, el empleo y el paro.
\item \emph{En la balanza de pagos}, los fondos estructurales y de cohesión se registran como \emph{transferencias de capital} (financian inversión) y los del FSE+ en la \emph{renta secundaria} (financian gasto corriente) \vertema{3.B.11}.
\end{itemize}

\begin{itemize}
\item \emph{Recapitulación}: la Unión converge, pero a 2 velocidades; los fondos tienen un efecto positivo, local y condicionado por la capacidad de cada región, que explica que el mismo dinero haya dado resultados tan distintos en el Este y en el Sur.
\end{itemize}

% ==========================================================================
\section*{Conclusión}
\minutos{1}

\begin{itemize}
\item \textbf{\emph{Recapitulación}}: En conclusión,
  \begin{itemize}
  \item La política regional existe porque la integración no reparte por igual: al abaratar el comercio, el mercado único puede concentrar la actividad en el centro. Se ejecuta con unos 392.000 millones en 2021-2027, en gestión compartida con los Estados y las regiones.
  \item En lo social y el empleo, la Unión pone mínimos, objetivos y coordinación, pero los Estados se reservan los salarios y la protección social.
  \item La convergencia real ha sido intensa en el Este y nula en el Sur desde 2008; los fondos ayudan, sobre todo donde hay capacidad para usarlos.
  \end{itemize}
\item \textbf{\emph{Relevancia}}:
  \begin{itemize}
  \item Para España, con 36.700 millones de euros en 2021-2027 y una renta que sigue por debajo del 92~\% de la media, la clave no es cuánto recibe sino en qué lo gasta \vertema{4.A.17} \vertema{4.A.30}.
  \end{itemize}
\item \textbf{\emph{Extensiones}}:
  \begin{itemize}
  \item El presupuesto que la financia \vertema{3.B.37}, la integración y la nueva geografía económica \vertema{3.B.6} \vertema{3.B.10}, las teorías y la evidencia del crecimiento \vertema{3.A.43} \vertema{3.A.45}, el federalismo fiscal \vertema{4.B.14} y el mercado de trabajo español \vertema{4.A.16}.
  \end{itemize}
\item \textbf{\emph{Opinión}}:
  \begin{itemize}
  \item La cohesión funciona mejor como política de inversión en capacidades (capital humano, instituciones, innovación) que como compensación o como caja para cada nueva prioridad. Simplificarla en planes nacionales puede ayudar a gastar, pero solo si no deja fuera a las regiones, que son quienes conocen sus problemas.
  \end{itemize}
\item \textbf{\emph{Idea final}}:
  \begin{itemize}
  \item Volvamos al refrán: Dios da pan a quien no tiene dientes. La evidencia dice que el pan sí alimenta y la política de cohesión ayuda a la convergencia de las regiones más desfavorecidas.
  \end{itemize}
\end{itemize}

\clearpage
\appendix

% ==========================================================================
\section{Anexo: el modelo de capital móvil:\cortetitulo{} integración, localización y cohesión}\label{anexo:capital-movil}


\paragraph{a) Idea}

\begin{itemize}
\item La riqueza, la población y el empleo están muy concentrados en el espacio: el «pentágono» entre Londres, París, Milán, Múnich y Hamburgo, o las capitales del Este. Esa concentración tiene determinantes de \emph{primera naturaleza} (recursos, clima, puertos) y de \emph{segunda naturaleza}, que son económicos: las empresas quieren estar cerca de sus clientes y de sus proveedores.
\item La \concepto{nueva geografía económica} de \textcite{Krugman1991Increasing}\footnote{\autor{Paul Krugman} recibió el Premio Nobel de Economía en 2008 «por su análisis de los patrones del comercio y de la localización de la actividad económica»\nocite{Nobel2008Krugman}.} explica cómo una pequeña asimetría puede convertirse en una gran diferencia entre regiones \vertema{3.B.6} \vertema{3.B.10}. Aquí usamos su versión más sencilla, el modelo de \emph{capital móvil} (\emph{footloose capital}) de \textcite{MartinRogers1995}, porque da una fórmula cerrada y sirve para comparar políticas regionales\nocite{BaldwinEtAl2003Economic}.
\end{itemize}

\paragraph{b) Modelo}

\subparagraph{i) Supuestos}

\begin{enumerate}
\item 2 regiones, el centro $N$ y la periferia $S$, con la misma tecnología. El centro es algo más grande: su cuota del gasto total es $s_E>1/2$.
\item Un sector agrícola con rendimientos constantes y comercio libre, que fija el salario, y un sector manufacturero con variedades diferenciadas de \autor{Dixit} y \autor{Stiglitz}, con elasticidad de sustitución $\sigma>1$ \vertema{3.A.18}.
\item Cada empresa manufacturera necesita una unidad de capital (su coste fijo). El \emph{capital se mueve} a la región donde rinde más, pero sus dueños no: gastan sus rentas donde viven. Por eso $s_E$ no cambia cuando se mueven las empresas.
\item Vender en la otra región cuesta $\tau>1$: de cada $\tau$ unidades enviadas llega una (coste iceberg). La \concepto{libertad comercial} es $\phi=\tau^{1-\sigma}$, que va de 0 (autarquía) a 1 (comercio libre). La integración sube $\phi$.
\end{enumerate}

\subparagraph{ii) Desarrollo}

\begin{itemize}
\item \emph{Paso 1: el beneficio en cada región}.
  \begin{itemize}
  \item Una empresa en $N$ vende en su mercado, el $s_E$ del gasto, y en el otro con el coste de $\tau$. Su beneficio es proporcional a $\frac{s_E}{\Delta_N}+\phi\frac{1-s_E}{\Delta_S}$, donde $\Delta_N=s_n+\phi(1-s_n)$ y $\Delta_S=\phi s_n+(1-s_n)$ miden la competencia que encuentra en cada mercado y $s_n$ es la cuota de empresas del centro.
  \item Más gasto local atrae empresas; más empresas locales, por la competencia, las repelen.
  \end{itemize}
\item \emph{Paso 2: el equilibrio de localización}. El capital se mueve hasta que el beneficio es igual en las 2 regiones. Igualando y despejando:
  \funcionclave[eq:mercado-domestico]{s_n=\frac12+\frac{1+\phi}{1-\phi}\left(s_E-\frac12\right)}
  \begin{itemize}
  \item \emph{Cómo se llega}:
    \begin{itemize}
    \item \emph{Demanda y precios}: con preferencias de \autor{Dixit} y \autor{Stiglitz}, cada empresa fija un precio igual a un margen fijo, $\frac{\sigma}{\sigma-1}$, sobre su coste marginal; como el salario es igual en las 2 regiones (lo fija la agricultura con comercio libre), todas las empresas cobran lo mismo en su mercado y $\tau$ veces más en el otro\nocite{BaldwinEtAl2003Economic}.
    \item \emph{Beneficio}: el beneficio de explotación de una empresa es la parte $1/\sigma$ de sus ventas. Con $K$ empresas (una por unidad de capital), gasto total $E$ y una parte $\mu$ de él en manufacturas:
      \[
      \pi_N=\frac{\mu E}{\sigma K}\left[\frac{s_E}{\Delta_N}+\phi\frac{1-s_E}{\Delta_S}\right],\qquad
      \pi_S=\frac{\mu E}{\sigma K}\left[\phi\frac{s_E}{\Delta_N}+\frac{1-s_E}{\Delta_S}\right],
      \]
      con $\Delta_N=s_n+\phi(1-s_n)$ y $\Delta_S=\phi s_n+(1-s_n)$.
    \item \emph{Equilibrio}: con $\pi_N=\pi_S$,
      \[
      (1-\phi)\frac{s_E}{\Delta_N}=(1-\phi)\frac{1-s_E}{\Delta_S}\;\Longrightarrow\; s_E\,\Delta_S=(1-s_E)\,\Delta_N .
      \]
      Sustituyendo $\Delta_N$ y $\Delta_S$ y despejando $s_n$ se llega a \eqref{eq:mercado-domestico}: $s_n=\frac12+\frac{1+\phi}{1-\phi}\left(s_E-\frac12\right)$, válida mientras $0\le s_n\le1$.\findemo
    \end{itemize}
  \end{itemize}
\item \emph{Paso 3: el efecto del mercado doméstico amplificado}.
  \begin{itemize}
  \item Como $\frac{1+\phi}{1-\phi}>1$, la cuota de empresas del centro es mayor que su cuota del gasto: el mercado grande atrae más que proporcionalmente a la industria.
  \item El factor crece con $\phi$: cuanto más integrado está el mercado, más fácil es servir a la periferia desde el centro.
  \item Con $s_n=1$ se llega a la aglomeración total, en $\phi^{\mathrm{CP}}=\frac{1-s_E}{s_E}$. Con $s_E=0{,}6$, en $\phi=2/3$ (gráfico~\ref{fig:mercado-domestico}).
  \end{itemize}
\end{itemize}

\begin{figure}[H]
\centering
\caption{El efecto del mercado doméstico amplificado: integración y transferencias}\label{fig:mercado-domestico}
\grafico[0.7\textwidth]{mercado-domestico}
\fuente{elaboración propia a partir de \textcite{MartinRogers1995} y \textcite{BaldwinEtAl2003Economic}, con $s_E=0{,}6$ y, tras la transferencia, $s_E'=0{,}55$.}
\end{figure}

\subparagraph{iii) Implicaciones}

\begin{itemize}
\item \emph{La integración no es neutral en el espacio}: al subir $\phi$, el mercado único lleva industria del Sur y del Este al centro, aunque eleve la renta del conjunto. Es la razón económica de la cohesión como compensación.
\item \emph{Las transferencias de renta a la periferia frenan la concentración}: si los fondos suben el gasto de $S$ (bajan $s_E$ a $s_E'$), la curva baja y la aglomeración total exige más integración ($\phi^{\mathrm{CP}\prime}>\phi^{\mathrm{CP}}$).
\item \emph{No toda infraestructura ayuda a la periferia} \nocite{MartinRogers1995,Martin1999}:
  \begin{itemize}
  \item la que une la periferia con el centro (una autovía o un tren rápido entre ambas) sube $\phi$ y puede vaciar la periferia, porque deja servirla desde el centro: es la «carretera de doble sentido»;
  \item la que abarata el comercio \emph{dentro} de la periferia agranda su mercado efectivo y atrae empresas.
  \end{itemize}
\item \emph{El dilema de la cohesión}: \textcite{Martin1999} añade crecimiento endógeno con efectos desbordamiento locales. Concentrar la innovación en el centro hace crecer más a toda la Unión, pero agranda la desigualdad entre regiones. Hay que elegir entre crecimiento y equidad espacial, salvo con políticas que abaraten la innovación y difundan el conocimiento, que mejoran las 2 cosas.
\end{itemize}

\paragraph{c) Extensiones}

\begin{itemize}
\item \emph{Por qué no hay causación circular}: como los dueños del capital gastan donde viven, mover una empresa no mueve gasto, y $s_E$ no depende de $s_n$. En el modelo centro-periferia los que se mueven son trabajadores, que se llevan su gasto, y la aglomeración se refuerza a sí misma \vertema{3.B.6}.
\item \emph{Causación circular}: en el modelo centro-periferia de \textcite{Krugman1991Increasing}, los trabajadores también se mueven y llevan su gasto, de modo que $s_E$ se vuelve endógena: aparecen puntos de ruptura, aglomeración total e histéresis (el diagrama \emph{tomahawk}) \vertema{3.B.6}.
\item \emph{La U de \textcite{KrugmanVenables1995Globalization}}: con vínculos entre industrias y salarios que reaccionan, la desigualdad entre centro y periferia primero sube y luego baja al integrarse. Con comercio muy barato, los salarios bajos de la periferia la vuelven a hacer atractiva: es lo que ha pasado con las fábricas que se han ido al Este.
\item \emph{Los efectos desbordamiento del conocimiento}: \textcite{BaldwinForslid2000Core} muestran que abaratar el comercio de bienes favorece la aglomeración, y abaratar el de ideas, la dispersión.
\end{itemize}

\paragraph{d) Evidencia empírica}

\begin{itemize}
\item \textcite{Puga2002} resume la de las primeras décadas: pese a los fondos, las diferencias entre \emph{países} se redujeron, y las diferencias entre \emph{regiones} de un mismo país crecieron, con más especialización y un paro más polarizado.
\item Con los datos de 2000-2024, la dispersión entre Estados sigue bajando y la dispersión entre regiones baja mucho menos desde 2008 (gráfico~\ref{fig:sigma}): las capitales del Este, como Varsovia (del 105~\% al 157~\% de la media) o Bucarest, despegan mientras sus regiones rurales se quedan atrás\nocite{Eurostat2026Regional}.
\item Las regiones de renta media que se estancan, con poca industria y poco capital humano, concentran el voto contrario a la Unión\nocite{IammarinoEtAl2019,RodriguezPose2018}.
\end{itemize}

\paragraph{e) Valoración}

\begin{itemize}
\item El modelo da la razón económica de que la cohesión acompañe al mercado único y al euro: la integración tiende a concentrar, y las transferencias pueden contrarrestarlo.
\item Advierte además contra la política regional «de carreteras»: unir mejor la periferia con el centro puede perjudicarla. Las prioridades actuales (innovación, capital humano, especialización inteligente) siguen esa lección.
\item Pero es un modelo de 2 regiones, sin congestión ni costes de vivienda en el centro, y no dice qué política funciona mejor en cada caso: eso es empírico (\ref{sec:papel}).
\end{itemize}

% ==========================================================================
\section{Anexo: los objetivos de la cohesión por periodos}\label{anexo:objetivos}

\begin{table}[H]
\centering
\caption{Objetivos y fondos de la política de cohesión, 1989-2027}\label{tab:objetivos}
\small
\begin{tabular}{>{\raggedright\arraybackslash}m{2.1cm}>{\raggedright\arraybackslash}m{7.6cm}>{\raggedright\arraybackslash}m{5cm}}
\hline
\textbf{Periodo} & \textbf{Objetivos} & \textbf{Rasgos}\\
\hline
1989-1993 & 5 objetivos: 1, regiones atrasadas (<75~\%); 2, reconversión industrial; 3 y 4, paro de larga duración y jóvenes; 5, agricultura y zonas rurales & reforma de 1988: concentración, programación, asociación y adicionalidad\\
1994-1999 & los mismos, más el 6 (zonas árticas, 1995) & Fondo de Cohesión (1993); Comité de las Regiones\\
2000-2006 & 3 objetivos: 1, regiones atrasadas; 2, reconversión; 3, empleo y formación & preadhesión para el Este; ampliación de 2004\\
2007-2013 & convergencia; competitividad regional y empleo; cooperación territorial & todas las regiones, con régimen transitorio; Estrategia de Lisboa\\
2014-2020 & 11 objetivos temáticos; 3 categorías de regiones (<75, 75-90, >90~\%) & Fondos Estructurales y de Inversión Europeos (con FEADER y pesca); condiciones \emph{ex ante}; REACT-EU (2020)\\
2021-2027 & 5 objetivos políticos y el del FTJ; 3 categorías (<75, 75-100, >100~\%) & RDC común a 8 fondos; condiciones favorecedoras; STEP (2024) y revisión intermedia (2025)\\
\hline
\end{tabular}
\fuente{elaboración propia a partir de la Comisión Europea\nocite{CE2018Panorama64,UE2021RDC} y del Parlamento Europeo\nocite{PE2026FichaCohesion}.}
\end{table}

% ==========================================================================
\section{Anexo: el reparto de competencias\cortetitulo{} de la Unión}\label{anexo:competencias}

\begin{itemize}
\item \emph{La regla de partida es la atribución}: la Unión solo tiene las competencias que le dan los Tratados, y las demás siguen siendo de los Estados (arts.~4.1 y 5 TUE). El art.~2 TFUE las ordena en categorías según quién legisla, y los arts.~3 a 6 enumeran los ámbitos de cada una (figura~\ref{figu:competencias-ue}). Las ejerce con \emph{proporcionalidad} (no más de lo necesario) y, fuera de las exclusivas, con \emph{subsidiariedad} (solo si el objetivo se alcanza mejor en común).
\item \emph{Exclusivas y compartidas}: en las \concepto{exclusivas} (art.~3) solo legisla la Unión; los Estados, solo si ella los faculta o para aplicar sus actos. En las \concepto{compartidas} (art.~4), la categoría residual, legislan las 2, pero lo que regula un acto de la Unión deja de estar disponible para los Estados: es la \emph{preferencia} u ocupación del campo (art.~2.2 y Protocolo n.º~25), que solo alcanza a lo que el acto regula. En investigación, espacio, cooperación para el desarrollo y ayuda humanitaria son \emph{paralelas}: la acción de la Unión no impide la de los Estados (art.~4.3 y 4.4).
\item \emph{Coordinación y apoyo}: en las políticas económicas y de empleo (art.~5) la política sigue siendo de los Estados y la Unión la coordina con orientaciones del Consejo; en los ámbitos del art.~6 (salud, industria, cultura, turismo, educación, protección civil, cooperación administrativa) la Unión solo apoya, coordina o complementa, sin armonizar las leyes nacionales. La PESC tiene un régimen propio, intergubernamental, con unanimidad y sin actos legislativos (art.~24 TUE).
\item \emph{Dónde está este tema}: la cohesión económica, social y territorial (art.~4.2.c; arts.~174 a 178) y la política social, en los aspectos que define el Tratado (art.~4.2.b; art.~153), son compartidas; las políticas de empleo (art.~5.2; arts.~145 a 150) y la coordinación de las sociales (art.~5.3) son de coordinación.
\item \emph{Qué se deduce}: por eso las 2 mitades del tema funcionan distinto. La cohesión se hace con dinero, en gestión compartida; en lo social, la Unión solo fija mínimos (y no en salarios, derecho de asociación, huelga ni cierre patronal, art.~153.5), y en el empleo solo orienta y vigila a través del Semestre Europeo.
\end{itemize}

\begin{infografia}[H]
\centering
\caption{El reparto de competencias de la Unión (arts.~2 a 6 TFUE): la cohesión, la política social y el empleo}\label{figu:competencias-ue}
\grafico[\textwidth]{competencias-ue}
\fuente{elaboración propia a partir de los arts.~2 a 6 del TFUE y de los arts.~4, 5 y 24 del TUE\nocite{UE2016TFUE}.}
\end{infografia}

% ==========================================================================
\section{Anexo: preguntas de test}\label{anexo:test}

\begin{anotaciones}
Preguntas tipo test del primer ejercicio: \href{https://www.victorgutierrezmarcos.es/oposicion/temario/primer-ejercicio/test/simulador.html?tema=3.B.40}{simulador de test}.
\end{anotaciones}

\bibliografia
\end{document}
